%% ma: L10-B14-0020
%% lop: 10
%% chuong: 5
%% bai: 14
%% dang: 10.14.1
%% loai: TN4
%% muc: NB
%% dap_an: "C"
%% nguon: "chuyen-de-cac-so-dac-trung-cua-mau-so-lieu-khong-ghep-nhom-toan-10-chuong-trinh-moi.pdf (GV. Phan Nhật Linh) (P04), Bài 03 Dạng 1 Phần I Câu 10"
%% kien_thuc: "Lớp 10 Bài 14 (khoảng biến thiên)"
%% trang_thai: chinh-thuc
%% ghi_chu: "Câu 13 gốc trùng hoàn toàn câu này, đã bỏ. Tự duyệt theo yêu cầu giáo viên 22/10/2026."
\begin{ex}
Mẫu số liệu sau cho biết chiều cao (đơn vị: cm) của các bạn học sinh trong tổ:
\begin{center}
$164$\quad $159$\quad $170$\quad $166$\quad $163$\quad $168$\quad $170$\quad $158$\quad $162$
\end{center}
Khoảng biến thiên $R$ của mẫu số liệu là
\choice{$R=10$}{$R=11$}{\True $R=12$}{$R=9$}
\loigiai{$R=170-158=12$. Chọn C.}
\end{ex}
